% =====================================================================
% LongRAG Full-Track Report
% A Controlled-Adaptation Baseline in a Multi-Architecture RAG Study
% =====================================================================
\documentclass[11pt,a4paper]{article}

\usepackage[utf8]{inputenc}
\usepackage[T1]{fontenc}
\usepackage{lmodern}
\usepackage[margin=1in]{geometry}
\usepackage{amsmath,amssymb,amsfonts}
\usepackage{booktabs}
\usepackage{multirow}
\usepackage{longtable}
\usepackage{array}
\usepackage{tabularx}
\usepackage{siunitx}
\usepackage{graphicx}
\usepackage{xcolor}
\usepackage{hyperref}
\usepackage{url}
\usepackage{enumitem}
\usepackage{listings}
\usepackage{fancyhdr}
\usepackage{titlesec}
\usepackage{caption}
\usepackage{subcaption}
\usepackage{footmisc}
\usepackage{microtype}
\usepackage{parskip}

\hypersetup{
  colorlinks=true,
  linkcolor=blue!50!black,
  citecolor=blue!50!black,
  urlcolor=blue!50!black,
  pdftitle={LongRAG Full-Track Report},
  pdfauthor={RAG Architecture Readiness Program},
}

\sisetup{
  round-mode=places,
  round-precision=4,
  detect-weight=true,
  detect-family=true,
}

\definecolor{codebg}{RGB}{248,248,248}
\definecolor{codekw}{RGB}{0,90,160}
\definecolor{codestr}{RGB}{160,50,20}
\definecolor{codecom}{RGB}{110,110,110}

\lstdefinelanguage{yaml}{
  keywords={true,false,null},
  keywordstyle=\color{codekw}\bfseries,
  commentstyle=\color{codecom}\itshape,
  stringstyle=\color{codestr},
  morecomment=[l]{\#},
}
\lstset{
  basicstyle=\ttfamily\small,
  backgroundcolor=\color{codebg},
  frame=single,
  rulecolor=\color{black!20},
  breaklines=true,
  showstringspaces=false,
  captionpos=b,
  columns=fullflexible,
  keepspaces=true,
}

\titleformat{\section}{\large\bfseries}{\thesection}{1em}{}
\titleformat{\subsection}{\normalsize\bfseries}{\thesubsection}{1em}{}

\pagestyle{fancy}
\fancyhf{}
\rhead{\small LongRAG Full-Track Report}
\lhead{\small Controlled-Adaptation Baseline}
\cfoot{\thepage}

\title{
  \vspace{-1.5cm}
  \rule{\linewidth}{0.4pt}\\[0.3cm]
  {\Large\bfseries LongRAG as a Controlled-Adaptation Baseline:}\\[0.15cm]
  {\large A Full-Pipeline Retrieval-Only Evaluation on Four Heterogeneous\\ Long-Context Corpora}\\[0.2cm]
  \rule{\linewidth}{0.4pt}\\[0.4cm]
  \normalsize RAG Architecture Readiness Program\\
  \small Technical Report --- LongRAG Track Closure
}
\author{}
\date{\today}

\begin{document}

\maketitle
\thispagestyle{empty}

% =====================================================================
\begin{abstract}
\noindent
This report documents the design, execution, and closure of a
full-pipeline retrieval-only evaluation of \emph{LongRAG}, a
long-context retrieval architecture originally proposed by
Jiang et al.\ (TIGER-AI-Lab, arXiv:2406.15319). The evaluation is
conducted as the \emph{first baseline} within a broader
seven-architecture comparative study whose goal is to answer the
question: \emph{which architecture retrieves the most relevant
information from long-context corpora under a shared, controlled,
retrieval-only protocol?} The architecture is treated as a
\textbf{controlled adaptation}: the architecture-defining components
(long-document grouping and dense-only retrieval) are preserved
exactly as described by the authors; the embedder is replaced by a
single pinned BGE-M3 (dense-only, revision
\texttt{5617a9f6\dots}) applied uniformly across all baselines in the
study; storage is provided by Qdrant with demonstrated
semantically-neutral equivalence on a frozen synthetic fixture. All
evaluation is conducted on four heterogeneous canonical datasets
(CLERC, CUAD, GreekBarBench, LaRA) drawn from a registered suite of
seven; the remaining three (RULER, Loong, LongBench-v2) were excluded
by explicit scoping decision, documented as a limitation. The
full pipeline consists of seven stages (L1--L7) with checksum-closed
artifacts, deterministic serialization, resume-safe execution, and a
strict isolation between retrieval and evaluation.
The primary result is
$\mathrm{macro\_document\_Recall@10} = \num{0.3144}$,
with strong per-dataset heterogeneity: LaRA (\num{0.8861}) dominates,
GreekBarBench is intermediate (\num{0.2606}), and CLERC
(\num{0.0786}) and CUAD (\num{0.0324}) are weak. We attribute the
weak performance on CLERC, CUAD, and GreekBarBench to the intrinsic
granularity of the paper's grouping policy---which packs $K$
consecutive documents into a single long unit of approximately
\SI{4096}{} tokens---when the corpus does not exhibit Wikipedia-style
hyperlink adjacency. This is reported as an honest architectural
finding, not as a defect. The final architecture status is
\textsc{LongRAG\_Full\_Pipeline\_Closed\_With\_Classified\_Limitations}.
\end{abstract}

\vfill
\noindent\textbf{Keywords.} retrieval-augmented generation;
long-context retrieval; controlled baseline comparison;
retrieval-only evaluation; reproducibility; academic benchmarking.

\newpage
\tableofcontents
\newpage

% =====================================================================
\section{Introduction}\label{sec:intro}

\subsection{Motivation}
Retrieval-Augmented Generation (RAG) has become the dominant
paradigm for grounding large language models on private and
long-tail corpora. In the long-context regime---where a single query
must be resolved against a corpus whose units span tens of
thousands to millions of tokens---a large and increasingly divergent
family of retrieval architectures has emerged. These range from
classical dense passage retrieval, through hybrid and learned-sparse
variants, to graph-based, hierarchical, and unit-aggregation
designs. There is currently no consensus on which design retrieves
relevant evidence most effectively when the corpus is long, when
documents are heterogeneous, and when the query set includes both
small and large contexts.

This report is part of a systematic effort to answer that question
with academic rigour. The broader program compares seven
architectures under a shared controlled protocol, with the following
commitments:

\begin{itemize}[leftmargin=1.5em,itemsep=0.15em]
  \item \textbf{Single embedder.} One frozen BGE-M3 (dense-only),
    same revision, across all architectures.
  \item \textbf{Single evaluator.} One canonical evaluator, one
    ground-truth projection, one metric set.
  \item \textbf{Retrieval-only.} No generation, no LLM-as-judge,
    no answer synthesis.
  \item \textbf{Architecture-defining preservation.} Every
    architecture's defining retrieval logic is preserved
    unchanged.
  \item \textbf{Documented substitutions.} Any replacement of a
    supporting component is documented, justified, and (where
    applicable) accompanied by an equivalence test.
\end{itemize}

LongRAG is the first architecture executed under this protocol. The
present document is its full closure report.

\subsection{Scope of this document}
This report covers, in order:
\begin{enumerate}[leftmargin=1.5em,itemsep=0.15em]
  \item the common controlled evaluation framework
    (\S\ref{sec:framework});
  \item the LongRAG architecture as defined by its authors and
    the exact fidelity boundary we adopt
    (\S\ref{sec:longrag});
  \item the seven-stage pipeline (L1--L7), including the
    \emph{fidelity correction} we performed mid-execution
    (\S\ref{sec:pipeline});
  \item the experimental setup: datasets, hardware, model, and
    measurement protocol (\S\ref{sec:experiments});
  \item results, per-dataset and macro-averaged
    (\S\ref{sec:results});
  \item a discussion of the observed behaviour, with emphasis on
    the role of document-level granularity
    (\S\ref{sec:discussion});
  \item limitations and threats to validity
    (\S\ref{sec:limitations});
  \item reproducibility artifacts, checksum chain, and closure
    manifest (\S\ref{sec:repro});
  \item appendices with full manifest excerpts and per-query
    diagnostic tables.
\end{enumerate}

\subsection{What this report is \emph{not}}
It is not a paper-faithful reproduction of the original LongRAG
experiments. The original paper evaluates LongRAG on a specific
multi-hop QA benchmark with the Wikipedia corpus, using a
proprietary long-context reader and a natural-language answer
evaluator. Our study is retrieval-only, uses no reader, and
operates on a heterogeneous canonical corpus constructed from
non-Wikipedia sources. Our results therefore describe the
\emph{controlled-adaptation behaviour} of LongRAG, not the paper's
headline numbers. Where the adaptation is non-trivial, we say so
explicitly. This is, we believe, the correct scientific framing for
a controlled baseline comparison.

% =====================================================================
\section{The Common Controlled Evaluation Framework}\label{sec:framework}

The seven-architecture comparison is governed by a master protocol
in which each architecture is subject to the same conditions.
We summarize the framework here; details specific to LongRAG follow
in \S\ref{sec:longrag}.

\subsection{Two distinct tracks}
Every architecture is executed along two independent tracks:

\begin{itemize}[leftmargin=1.5em,itemsep=0.15em]
  \item \textbf{Controlled architectural comparison (primary).}
    Shared dataset envelope, shared embedder, shared evaluator,
    shared query set, shared ground-truth projection. The
    architecture-defining retrieval logic is preserved; every other
    component is either preserved, replaced with a documented
    equivalent, or removed if downstream.
  \item \textbf{Native-fidelity reference (secondary, where
    feasible).} Uses each architecture's original embedder,
    original storage, original hyperparameters. Its purpose is
    diagnostic: to detect whether the controlled adaptation changed
    the architecture's behaviour materially. Native fidelity is
    required to be small; it need not run on the full corpus.
\end{itemize}

For LongRAG specifically, the \emph{native track is not eligible}
because the architecture's original embedder identity is unpinned
in the reference repository (it depends on a Wikipedia-specific
encoder whose weights are not reproducibly versioned at the
revision level). This is recorded as a formal declaration in the
readiness contract, and no native-equivalence claim is made
anywhere in this report.

\subsection{Component fidelity boundary}
Every component of LongRAG is classified into exactly one of:

\begin{center}
\begin{tabular}{p{4.1cm} p{10.5cm}}
\toprule
\textbf{Category} & \textbf{Treatment} \\
\midrule
\textsc{Architecture Defining} &
  Preserved exactly. If the component cannot be executed under
  the controlled environment, the architecture is deferred with a
  reason; it is never approximated. \\
\textsc{Architecture Supporting} &
  Replaced only under a documented semantic-equivalence
  justification. Every replacement is accompanied by a
  \emph{substitution class} and, where possible, a
  \emph{top-$k$ equivalence test}. \\
\textsc{Optional Downstream} &
  Disabled for the retrieval-only benchmark. This includes
  generation, answer synthesis, LLM-as-judge, and any
  query-time reader. \\
\bottomrule
\end{tabular}
\end{center}

\subsection{Retrieval-only boundary}
The entire pipeline is retrieval-only. No generation is performed.
No answer is synthesized. No LLM judge is invoked. The
\emph{only} place ground truth is loaded is inside the evaluation
stage (L6), which runs as an independent process from the
retriever. This isolation is enforced by architecture: the retriever
receives only \texttt{(query\_id, query\_text)} pairs; the
evaluator receives \texttt{(query\_id, ranked\_doc\_ids)} plus the
canonical \texttt{(query\_id, relevant\_doc\_ids)} pairs from the
dataset.

\subsection{Common output contract}
Every retrieved query is serialized as:
\begin{lstlisting}[language=Python]
{
  "query_id": <str>,
  "ranked_items": [
    {
      "architecture_item_id":  <str>,
      "canonical_doc_id":      <str>,
      "canonical_chunk_or_evidence_ids": [<str>, ...],
      "rank":                  <int, >= 1, contiguous>,
      "score":                 <finite float>,
      "score_semantics":       "HIGHER_IS_BETTER",
      "metadata":              {<str>: <any>, ...}
    },
    ...
  ],
  "metadata": {<str>: <any>, ...}
}
\end{lstlisting}

The comparison between architectures is performed exclusively at
this interface. Storage backends are never compared directly.
This guarantees that architectural differences in retrieval logic
are the only differences visible to the evaluator.

% =====================================================================
\section{LongRAG: The Architecture, As Defined by Its Authors}\label{sec:longrag}

\subsection{Original architecture}
LongRAG\footnote{Jiang, Z.\ et al.\ \emph{LongRAG: Enhancing
Retrieval-Augmented Generation with Long-context LLMs},
arXiv:2406.15319, TIGER-AI-Lab.} is a retrieval architecture
designed for the long-context regime. Its central idea is to
\emph{reduce the corpus size} by grouping multiple source documents
into \emph{long retrieval units} of roughly fixed size, and to
retrieve at that coarser granularity. The motivating observation is
that in long-context tasks, the retrieval unit does not need to be
a short passage; it can be a long, self-contained bundle of
adjacent evidence.

The architecture is defined by two architectural commitments:

\begin{enumerate}[leftmargin=1.5em,itemsep=0.15em]
  \item \textbf{Long-document grouping.} $K$ consecutive source
    documents are packed into a single long unit whose token count
    is bounded above by an architectural target (typically
    $\approx \SI{4096}{}$ tokens). The choice of $K$ is not fixed;
    it is determined dynamically by the packing process until the
    token budget is filled. Grouping is \emph{adjacency-driven}:
    the paper's original setting exploits Wikipedia's hyperlink
    structure to define adjacency, but any deterministic
    adjacency policy that preserves source order is admissible.
  \item \textbf{Dense-only retrieval.} Retrieval is performed by
    computing dense vector similarity between a query vector and
    the long-unit vectors. No sparse signal, no hybrid fusion, no
    reranker, and no cross-encoder are used at retrieval time.
    Retrieval returns a small number of long units; the
    document-level ranking is obtained by expanding each retrieved
    unit to its member documents.
\end{enumerate}

LongRAG additionally includes a \emph{long reader} --- a
long-context LLM that consumes the retrieved long units and
synthesizes a final answer. Because this study is retrieval-only,
the long reader is disabled and is not part of our pipeline. This
is not a fidelity violation; it is the intended operation of the
architecture when the task is retrieval rather than generation.

\subsection{Why dense-only retrieval is architecture-defining}
A natural question is whether the dense-only constraint could be
relaxed by adding a sparse or hybrid signal. We argue it cannot,
for the following reason. The paper's primary claim is not that
dense retrieval is better than hybrid retrieval in general; it is
that \emph{the architectural reduction of the corpus to a small
number of long units} changes the retrieval regime. Once the
corpus is small (a few thousand long units rather than hundreds of
thousands of passages), the interaction between a hybrid signal and
the coarse unit granularity is not the same as for short-passage
retrieval. Introducing sparse or hybrid retrieval would therefore
change the architecture under test, not merely support it.

We therefore treat dense-only retrieval as \textsc{Architecture
Defining} and preserve it.

\subsection{Our controlled adaptation: exact choices}
\label{sec:adaptation}

The controlled adaptation of LongRAG used in this study is fixed by
four frozen contracts. We restate the resulting configuration here.

\begin{center}
\begin{tabular}{p{4.5cm} p{10cm}}
\toprule
\textbf{Property} & \textbf{Value / Decision} \\
\midrule
Fidelity class & \textsc{Controlled Adaptation} \\
Source architecture commit & \texttt{ba92f95} \\
Native track eligibility & Not eligible (embedder unpinned) \\
Long-document grouping policy & \texttt{adjacent\_document\_segments} \\
Long-unit token budget & \SI{4096}{} tokens \\
Token counter & Pinned BGE-M3 tokenizer,
  \texttt{add\_special\_tokens=False} \\
Oversize handling & Deterministic windowing of the offending
  document into \SI{4096}{}-token windows using the pinned
  BGE-M3 tokenizer \\
Retrieval signal & Dense-only cosine \\
Embedder & \texttt{BAAI/bge-m3} \\
Embedder revision & \texttt{5617a9f61b028005a4858fdac845db406aefb181} \\
Embedding dimension & 1024 \\
Embedding normalization & $\ell_2$, unit norm \\
Document embeddings precision & FP16 \\
Query embeddings precision & FP16 \\
Query embedding max length & \SI{4098}{} tokens \\
Storage backend & Qdrant (local persistence) \\
Storage vector name & \texttt{dense} \\
Storage distance & COSINE \\
Storage substitution class &
  \textsc{Semantically Neutral for Controlled Track}
  (see \S\ref{sec:storage-equiv}) \\
Retrieval depth & top-\num{100} long units \\
Document limit & \num{100} unique documents \\
Expansion policy & first occurrence by long-unit rank \\
Reader / generation / LLM judge & Disabled \\
\bottomrule
\end{tabular}
\end{center}

\subsection{Component disposition}
For each component we record its role, treatment, justification,
fidelity risk, and validation criterion. Table~\ref{tab:disposition}
is a faithful summary of the machine-readable contract stored at
\texttt{baseline\_audit/longrag\_full\_track\_component\_disposition.yaml}.

\begin{table}[h]
\centering
\caption{LongRAG component disposition under the controlled track.}
\label{tab:disposition}
\small
\begin{tabular}{p{3.1cm} p{3cm} p{7.9cm}}
\toprule
\textbf{Component} & \textbf{Role} & \textbf{Treatment and justification} \\
\midrule
\texttt{long\_document\_grouping} &
  \textsc{Arch Defining} &
  Preserved. Replacing with flat chunks would collapse LongRAG into a
  different architecture. \\
\texttt{dense\_only\_retrieval} &
  \textsc{Arch Defining} &
  Preserved. Introducing sparse / hybrid / reranker would change the
  architecture. \\
\texttt{canonical\_document\_input} &
  \textsc{Arch Supporting} &
  Controlled adaptation from Wikipedia-specific input to canonical
  documents. Low risk. \\
\texttt{embedder} &
  \textsc{Arch Supporting} &
  Controlled replacement to BGE-M3 dense-only, pinned revision.
  Equivalence with native embedder \emph{not claimed}. \\
\texttt{tokenizer} &
  \textsc{Arch Supporting} &
  Matched to embedder identity. No substitution. \\
\texttt{storage\_backend} &
  \textsc{Arch Supporting} &
  Qdrant dense-only, exact equivalence demonstrated on frozen fixture
  (\S\ref{sec:storage-equiv}). \\
\texttt{oversize\_handling} &
  \textsc{Arch Supporting} &
  Deterministic BGE-M3-tokenizer windowing for documents exceeding
  \SI{4096}{} tokens. Needed because canonical corpora are
  heterogeneous. \\
\texttt{tokenizer\_bound\_segment\_tokens} &
  \textsc{Arch Supporting} &
  Token counts measured with the same tokenizer that produces
  embeddings, so the \SI{4096}{}-token budget corresponds to the
  model input budget. \\
\texttt{llm\_index\_time / query\_time} &
  Not present &
  LongRAG has no LLM in retrieval. \\
\texttt{llm\_final\_generation} &
  \textsc{Optional Downstream} &
  Disabled. Retrieval-only benchmark. \\
\texttt{long\_reader} &
  \textsc{Optional Downstream} &
  Disabled. Explicitly excluded by the environment plan. \\
\bottomrule
\end{tabular}
\end{table}

% =====================================================================
\section{The Seven-Stage Pipeline}\label{sec:pipeline}

The pipeline is organized into seven named stages L1--L7. Each stage
produces at least one manifest with SHA-256 closure over all of its
outputs, is resume-safe (a stage marked \textsc{Succeeded} is not
re-executed unless \texttt{--force} is passed), and is atomic
(temporary file + \texttt{os.replace}, never partial writes).

\begin{center}
\begin{tabular}{c p{4.4cm} p{8cm}}
\toprule
\textbf{Stage} & \textbf{Purpose} & \textbf{Primary outputs} \\
\midrule
L1  & Architecture and dataset-envelope freeze &
  Frozen contracts, verification report, checksum manifest. \\
L1a & Mid-execution correction of the long-unit semantics &
  Contract version bump 0.1.0 $\to$ 0.2.0. \\
L1b & Fidelity correction of the grouping policy &
  Contract version bump 0.2.0 $\to$ 0.3.0; switches policy from
  \texttt{single\_document\_if\_within\_budget} to
  \texttt{adjacent\_document\_segments}. \\
L2  & Unit materialization &
  \texttt{units.jsonl} + per-dataset manifest and checksums. \\
L3  & Colab-GPU embeddings &
  Input/output bundles with SHA-256 closure; \texttt{embedding\_records.jsonl}
  per dataset. \\
L4  & Qdrant storage build and reopen verification &
  One collection per dataset; \texttt{storage\_manifest.json}. \\
L5  & Full retrieval &
  \texttt{results.jsonl} per query, with conformance to the common
  output contract. \\
L6  & Ground-truth evaluation &
  Per-query metrics, per-dataset aggregates, cross-dataset macro
  aggregates. \\
L7  & Closure &
  Cross-dataset registry, closure manifest, \LaTeX{} and Markdown
  reports. \\
\bottomrule
\end{tabular}
\end{center}

\subsection{L1 --- Architecture and dataset-envelope freeze}
The first stage loads the seven frozen LongRAG contracts, verifies
every canonical dataset's SHA-256 against the external canonical
suite registry, streams the canonical loaders to confirm document
and example counts, and writes:

\begin{itemize}[leftmargin=1.5em,itemsep=0.15em]
  \item \texttt{longrag\_full\_track\_embedding\_profile.yaml}
    (Colab GPU, FP16 primary; CPU FP32 bounded reference for
    equivalence only);
  \item \texttt{longrag\_full\_track\_architecture\_freeze\_contract.yaml};
  \item \texttt{longrag\_full\_track\_component\_disposition.yaml};
  \item \texttt{longrag\_full\_track\_dataset\_envelope.yaml};
  \item \texttt{longrag\_full\_track\_plan.yaml};
  \item a verification report and a checksum manifest.
\end{itemize}

\paragraph{Initial freeze (v0.1.0).}
The very first frozen contract set carried
\texttt{target\_token\_size = 256} and \texttt{construction\_policy =
adjacent\_segments}. This was informed by the synthetic acceptance
fixture (\S\ref{sec:synthetic}), which used twelve synthetic
documents and a small unit budget.

\subsection{L1a --- Correction of the long-unit semantics (v0.2.0)}
During reconnaissance it became clear that the architecture plan in
\texttt{longrag\_environment\_plan.yaml} specifies
\texttt{primary\_unit\_size\_tokens = 4096}, and the reference
implementation groups long units at the \SI{4096}{}-token scale, not
at \SI{256}{}. We therefore amended the freeze to
\texttt{target\_token\_size = 4096}, upgraded the token counter from
whitespace to the pinned BGE-M3 tokenizer, and introduced explicit
oversize handling. The contract version was bumped from
$0.1.0 \to 0.2.0$. This is a \emph{scale correction} and does not by
itself change the architecture's retrieval semantics.

\subsection{L1b --- Fidelity correction of the grouping policy (v0.3.0)}
\label{sec:fidelity-fix}

A more consequential correction was required shortly after. A careful
reading of the LongRAG paper shows that the architecture's defining
behaviour is to \emph{group $K$ consecutive documents into a single
long unit of $\approx \SI{4096}{}$ tokens}. Our v0.2.0 freeze had
instead adopted
\texttt{single\_document\_if\_within\_budget}, in which each source
document becomes its own long unit unless it exceeds the budget.

The difference is not cosmetic. Under
\texttt{single\_document\_if\_within\_budget}, a corpus of $N$ small
documents produces $N$ short units; the architectural claim that the
corpus is reduced to a small number of long units does not hold. On
CLERC, for example, v0.2.0 produced \num{2756} units with mean
length $\approx \SI{363}{}$ tokens, whereas the paper's grouping
policy produces \num{256} units with mean length $\approx
\SI{3900}{}$ tokens---a tenfold reduction of the retrieval corpus,
which is the central design property of LongRAG.

We therefore re-froze the architecture with:

\begin{itemize}[leftmargin=1.5em,itemsep=0.15em]
  \item \texttt{construction\_policy = adjacent\_document\_segments};
  \item \texttt{target\_token\_size = 4096};
  \item oversize handling via deterministic BGE-M3-tokenizer windows;
  \item contract version $0.2.0 \to 0.3.0$.
\end{itemize}

The prior freeze was retained in the git history; both are preserved
in the closure registry. We report this correction explicitly, in
the spirit of reproducible research, because omitting it would make
the reported numbers inconsistent with the paper's architecture.

\subsection{L2 --- Unit materialization}
For each dataset, L2 streams the canonical documents, builds one
source segment per document (with a deterministic
\texttt{parent\_id = doc\_id:p0} and \texttt{chunk\_id = seg-doc\_id}),
tokenizes each with the pinned BGE-M3 tokenizer, and constructs
long units under the frozen policy. The critical invariants are:

\begin{itemize}[leftmargin=1.5em,itemsep=0.15em]
  \item every unit has $\leq \SI{4096}{}$ BGE-M3 tokens;
  \item the union of unit \texttt{member\_doc\_ids} equals the
    set of source document ids;
  \item the union of unit \texttt{member\_chunk\_ids}, augmented by
    the \texttt{source\_segment\_chunk\_id} metadata field for split
    units, equals the set of source chunk ids;
  \item all unit ids are deterministic content addresses built from
    \texttt{(dataset\_id, policy, target\_token\_size, ordered
    chunk\_ids)}.
\end{itemize}

The materialization results are summarized in
Table~\ref{tab:l2}.

\begin{table}[h]
\centering
\caption{L2 unit materialization, per dataset. Mean unit length is
weighted by token count. Split count is the number of long units that
resulted from an oversize document being windowed.}
\label{tab:l2}
\small
\begin{tabular}{l r r r r r}
\toprule
\textbf{Dataset} & \textbf{Documents} & \textbf{Long units}
  & \textbf{Split units} & \textbf{Max tokens} & \textbf{Secs} \\
\midrule
CLERC              & \num{2756} & \num{256}   & \num{0}     & \num{4096} & \num{27.6} \\
CUAD               & \num{510}  & \num{1725}  & \num{1591}  & \num{4096} & \num{48.0} \\
GreekBarBench      & \num{284}  & \num{1437}  & \num{1427}  & \num{4096} & \num{45.4} \\
LaRA               & \num{131}  & \num{2251}  & \num{2251}  & \num{4096} & \num{59.1} \\
\midrule
\textbf{Total}     & \num{3681} & \num{5669}  & \num{5269}  & ---        & \num{180.1} \\
\bottomrule
\end{tabular}
\end{table}

\paragraph{Observation.}
The grouping policy produces genuine packing only on CLERC, where the
average document length ($\approx \SI{400}{}$ tokens) is an order
of magnitude smaller than the unit budget. On the other three
datasets, average document length exceeds the unit budget, so the
oversize-windowing path dominates: each source document becomes
multiple long units, and inter-document packing is negligible. This
observation is central to the interpretation of results
(\S\ref{sec:discussion}).

\subsection{L3 --- Embeddings (Colab GPU, FP16)}
\label{sec:l3}

Embeddings are produced exclusively on a Google Colab Tesla T4 GPU,
in FP16, using the pinned BGE-M3 model
(\texttt{revision 5617a9f6\dots}). The local workstation never loads
the model for document embedding. Each dataset is packed into a
self-contained input bundle that includes:

\begin{itemize}[leftmargin=1.5em,itemsep=0.15em]
  \item \texttt{units.jsonl};
  \item the L2 manifest, summary, and checksums;
  \item the Colab runner script;
  \item an input manifest recording SHA-256 of the unit file.
\end{itemize}

The Colab runner:

\begin{enumerate}[leftmargin=1.5em,itemsep=0.15em]
  \item verifies the input SHA-256 against the embedded manifest;
  \item loads BGE-M3 dense-only, FP16, CUDA;
  \item encodes all units, one batch at a time, with
    \texttt{return\_dense=True}, \texttt{return\_sparse=False},
    \texttt{return\_colbert\_vecs=False};
  \item verifies the number of encoded units equals the number of
    input units;
  \item writes \texttt{embedding\_records.jsonl} with one record per
    unit and \texttt{embedding\_manifest.json} with the config hash
    and the FP16 / GPU / max-length provenance;
  \item packages the output bundle.
\end{enumerate}

Upon completion, the output bundles are transferred back to the
local workstation where L3-import re-validates the schema using
Pydantic, cross-checks the input unit SHA-256 from the output
manifest, and verifies the record count. Any mismatch aborts the
import for that dataset only.

\subsection{L4 --- Qdrant storage and reopen equivalence}
The embeddings are written into a local Qdrant instance with one
collection per dataset
(\texttt{longrag-<dataset>-v1}). The write path is:

\begin{enumerate}[leftmargin=1.5em,itemsep=0.15em]
  \item create the collection with a single named dense vector of
    dimension \num{1024}, distance \textsc{Cosine};
  \item upsert points in idempotent batches;
  \item close the client;
  \item reopen a fresh client on the same directory;
  \item verify that the reopened collection's point count equals the
    expected count.
\end{enumerate}

Table~\ref{tab:l4} summarizes the L4 results. Every collection
reopened with the correct count; the ingest step is idempotent and
can be safely re-executed.

\begin{table}[h]
\centering
\caption{L4 Qdrant storage build and reopen verification.}
\label{tab:l4}
\small
\begin{tabular}{l r r c r}
\toprule
\textbf{Dataset} & \textbf{Expected} & \textbf{Ingested}
  & \textbf{Reopen} & \textbf{Secs} \\
\midrule
CLERC              & \num{256}   & \num{256}   & OK & \num{0.73} \\
CUAD               & \num{1725}  & \num{1725}  & OK & \num{4.89} \\
GreekBarBench      & \num{1437}  & \num{1437}  & OK & \num{4.04} \\
LaRA               & \num{2251}  & \num{2251}  & OK & \num{6.33} \\
\midrule
\textbf{Total}     & \num{5669}  & \num{5669}  & OK & \num{16.0} \\
\bottomrule
\end{tabular}
\end{table}

\subsubsection{Storage equivalence}
\label{sec:storage-equiv}
The storage substitution is classified as
\textsc{Semantically Neutral for Controlled Track}. This
classification was established by an earlier bounded synthetic
acceptance test whose report
(\texttt{longrag\_storage\_equivalence\_report.yaml}) records:

\begin{itemize}[leftmargin=1.5em,itemsep=0.15em]
  \item top-1, top-5, and top-10 overlap of \num{1.0} between Qdrant
    and an exact in-memory dense reference index, at every cutoff;
  \item maximum absolute score difference
    $\num{1.78e-15}$;
  \item mean absolute score difference $\num{6.29e-16}$;
  \item exact identity and order matching after tie normalization.
\end{itemize}

The equivalence is established \emph{only} on the small frozen
fixture. Large-corpus approximate-search settings require a separate
contract; no approximate HNSW override is applied in the Qdrant
collections used here.

\subsection{L5 --- Full retrieval}
L5 retrieves the top-\num{100} long units for each query and expands
them to a ranked list of unique canonical documents.

\paragraph{Query projection.}
For every example, only \texttt{(query\_id, question)} is retained.
The remaining fields---including \texttt{ground\_truth.relevant\_ids},
\texttt{ground\_truth.answers}, \texttt{ground\_truth.spans},
\texttt{ground\_truth.choice}, and \texttt{query.doc\_ids}---are
discarded before retrieval. The projection is a typed model whose
construction reads the canonical example object and emits a
\emph{lean} record. This implements the ground-truth isolation
required by the master protocol.

\paragraph{Query embeddings.}
Query embeddings are computed \emph{once} per dataset on the same
Colab GPU used for document embeddings, in FP16, with
\texttt{max\_length = 4098} (matching the document embedding
setting). The embeddings are cached to
\texttt{outputs/external\_longrag/\_query\_embeddings/}. This cache
is \emph{reusable}: query embeddings depend only on
\texttt{(query\_text, model, revision, precision, max\_length)} and
not on the storage or architecture under test. The same cache will
be consumed by subsequent architectures in the program.

\paragraph{Retrieval.}
For each query, we compute a cosine similarity search on the
dataset's Qdrant collection, retrieve the top-\num{100} long units,
convert them into a \texttt{RankedLongUnit} sequence, and expand to
canonical documents by iterating the retrieved units in rank order
and emitting each unseen document once. The expansion policy is
\texttt{first\_occurrence\_by\_long\_unit\_rank}, exactly as
implemented by the reference adapter in
\texttt{rag\_system.baselines.longrag}.

Table~\ref{tab:l5} reports the retrieval timing.

\begin{table}[h]
\centering
\caption{L5 retrieval timing. Query encoding is off the critical
path here: it was executed once on Colab GPU and cached. The
retrieval time shown is Qdrant retrieval + expansion only.}
\label{tab:l5}
\small
\begin{tabular}{l r r r}
\toprule
\textbf{Dataset} & \textbf{Queries} & \textbf{Retrieval (s)}
  & \textbf{Throughput (q/s)} \\
\midrule
CLERC              & \num{1000}  & \num{6.56}   & \num{152.5} \\
CUAD               & \num{20910} & \num{294.29} & \num{71.0} \\
GreekBarBench      & \num{284}   & \num{3.56}   & \num{79.8} \\
LaRA               & \num{2326}  & \num{37.56}  & \num{61.9} \\
\midrule
\textbf{Total}     & \num{24520} & \num{341.97} & --- \\
\bottomrule
\end{tabular}
\end{table}

\subsection{L6 --- Ground-truth evaluation}
L6 runs as an independent process from L5. It loads the retrieval
results and the canonical ground truth, verifies that the query-id
sets match, and computes the following metrics on canonical
document ids:

\begin{itemize}[leftmargin=1.5em,itemsep=0.15em]
  \item \textbf{Hit@$k$} for $k \in \{1, 5, 10, 20, 50\}$:
    1 if at least one relevant document appears within the first
    $k$ retrieved documents, else 0.
  \item \textbf{Recall@$k$} for the same cutoffs: the number of
    relevant documents appearing within the first $k$, divided by
    the total number of relevant documents for that query.
  \item \textbf{Precision@10}: relevant documents in the first ten
    positions divided by ten.
  \item \textbf{MRR@$k$} for the same cutoffs: reciprocal of the
    rank of the first relevant document within the cutoff, or zero
    if none.
  \item \textbf{nDCG@$k$} for $k \in \{10, 50\}$: gain formula
    $2^{\mathrm{rel}} - 1$ with discount $1/\log_2(i+1)$, normalized
    by the ideal DCG computed from the query's full relevant set.
\end{itemize}

\paragraph{Zero-relevance queries.}
Following the frozen evaluator contract
(\texttt{common\_retrieval\_evaluation\_contract.yaml}), queries
with no relevant documents are \emph{excluded} from all metric
aggregations, and the number of excluded queries is reported
explicitly per dataset. On CUAD this is significant:
\num{14208} of \num{20910} queries ($\approx 68\%$) have zero
relevant documents and are excluded.

\paragraph{Aggregation.}
For each dataset we compute a macro-average over evaluated queries.
Across datasets we compute a second-level macro-average with equal
weight per dataset. The primary metric is
$\mathrm{macro\_document\_Recall@10}$.

\subsection{L7 --- Closure}
L7 collects the state of every dataset, verifies that all stage
manifests exist, and writes:

\begin{itemize}[leftmargin=1.5em,itemsep=0.15em]
  \item \texttt{baseline\_audit/longrag\_full\_track\_closure\_manifest.yaml};
  \item \texttt{baseline\_audit/longrag\_full\_track\_cross\_dataset\_registry.yaml};
  \item \texttt{reports/longrag\_full\_track\_closure\_report.tex};
  \item \texttt{docs/longrag\_full\_track\_closure.md}.
\end{itemize}

The closure status is set to
\textsc{LongRAG\_Full\_Pipeline\_Closed\_With\_Classified\_Limitations}.

% =====================================================================
\section{Experimental Setup}\label{sec:experiments}

\subsection{Datasets}
The canonical suite contains seven datasets, of which four were
executed. Table~\ref{tab:datasets} summarizes them.

\begin{table}[h]
\centering
\caption{Executed datasets. The canonical suite contains three
additional datasets (RULER, Loong, LongBench-v2) that were
excluded by explicit scoping decision (see \S\ref{sec:limitations}).}
\label{tab:datasets}
\small
\begin{tabular}{l r r l}
\toprule
\textbf{Dataset} & \textbf{Documents} & \textbf{Examples}
  & \textbf{Scope} \\
\midrule
CLERC              & \num{2756} & \num{1000}
  & Global dataset corpus \\
CUAD               & \num{510}  & \num{20910}
  & Shared domain corpus \\
GreekBarBench      & \num{284}  & \num{284}
  & Query-specific corpus \\
LaRA               & \num{131}  & \num{2326}
  & Shared domain corpus \\
\midrule
\textbf{Total}     & \num{3681} & \num{24520} & \\
\bottomrule
\end{tabular}
\end{table}

For every dataset the canonical \texttt{documents.jsonl} and
\texttt{examples.jsonl} SHA-256 was verified against the external
canonical suite registry before any pipeline stage ran. The
registry fingerprint is
\texttt{5f577978b6ab7a4ee3fef29b68db29f95b37261f103bb6cfd5fd7094c70a2aa4}.

\subsection{Model and precision}
All dense vectors in this study are produced by BGE-M3
(\texttt{BAAI/bge-m3}) at revision
\texttt{5617a9f61b028005a4858fdac845db406aefb181}, in FP16, on a
Google Colab Tesla T4 GPU. Both document embeddings and query
embeddings use the same model, precision, and tokenizer. Document
embeddings use no input truncation, because long-unit token counts
are bounded above by \SI{4096}{} tokens at the L2 stage. Query
embeddings use \texttt{max\_length = 4098} (i.e.\ \SI{4096}{} content
tokens plus the two BGE-M3 special tokens).

The decision to use FP16 uniformly is a \emph{controlled
adaptation} of the original architecture. The original LongRAG paper
uses a proprietary encoder; no revision-level identity is publicly
pinned. FP16 BGE-M3 is applied identically to every architecture in
the study, so comparisons between architectures are unaffected by
this choice.

\subsection{Hardware}
\begin{itemize}[leftmargin=1.5em,itemsep=0.15em]
  \item \textbf{Embedding (L3, and L5 query encoding):}
    Google Colab, Tesla T4, 16\,GiB VRAM, CUDA runtime, FP16.
  \item \textbf{All other stages (L1, L2, L4, L5 retrieval, L6, L7):}
    Local workstation, Windows 11, PowerShell shell,
    Python 3.12.10, CPU only.
\end{itemize}

\subsection{Software stack}
\begin{center}
\begin{tabular}{l l}
\toprule
\textbf{Component} & \textbf{Version} \\
\midrule
Python               & 3.12.10 \\
FlagEmbedding        & 1.4.2 \\
\texttt{qdrant-client} & 1.19.1 \\
Pydantic             & (project-pinned) \\
Ruff                 & (project-configured, line length 100) \\
mypy                 & (project-configured, strict for \texttt{rag\_system}) \\
\bottomrule
\end{tabular}
\end{center}

\subsection{Synthetic acceptance test}
\label{sec:synthetic}
Prior to full execution, a bounded synthetic acceptance run was
executed on a twelve-document fixture with six queries. This
verified the end-to-end pipeline (canonical input,
long-unit construction, dense embedding, reference index, Qdrant
index, retrieval, expansion, evaluation case construction) in
under a minute, using the same code paths as the full run. The
acceptance contract is stored at
\texttt{baseline\_audit/longrag\_synthetic\_acceptance\_contract.yaml}
and its manifest at
\texttt{baseline\_audit/longrag\_synthetic\_acceptance\_manifest.yaml}.

% =====================================================================
\section{Results}\label{sec:results}

\subsection{Per-dataset metrics}
Table~\ref{tab:results} reports the complete per-dataset results. All
metrics are computed on canonical document ids. Zero-relevance
queries are excluded with explicit counts. Every number in this table
is reproduced verbatim from
\texttt{outputs/external\_longrag/\_l6\_cross\_dataset\_evaluation.json}.

\begin{table}[h]
\centering
\caption{Per-dataset and macro-averaged metrics. Evaluated counts
exclude zero-relevance queries; the number of excluded queries is
shown in column 3.}
\label{tab:results}
\small
\begin{tabular}{l r r r r r r r r}
\toprule
\textbf{Dataset} & \textbf{Ev.} & \textbf{Excl.}
  & \textbf{H@1} & \textbf{H@5} & \textbf{H@10} & \textbf{H@50}
  & \textbf{R@10} & \textbf{nDCG@10} \\
\midrule
CLERC              & \num{1000} & \num{0}
  & \num{0.0310} & \num{0.0550} & \num{0.0940} & \num{0.2110}
  & \num{0.0786} & \num{0.0543} \\
CUAD               & \num{6702} & \num{14208}
  & \num{0.0039} & \num{0.0161} & \num{0.0324} & \num{0.1458}
  & \num{0.0324} & \num{0.0153} \\
GreekBarBench      & \num{284}  & \num{0}
  & \num{0.0352} & \num{0.1303} & \num{0.2606} & \num{0.6338}
  & \num{0.2606} & \num{0.1233} \\
LaRA               & \num{2326} & \num{0}
  & \num{0.7128} & \num{0.8409} & \num{0.8861} & \num{0.9811}
  & \num{0.8861} & \num{0.7977} \\
\midrule
\textbf{Macro}     & --- & ---
  & \num{0.1957} & \num{0.2606} & \num{0.3183} & \num{0.4929}
  & \num{0.3144} & \num{0.2476} \\
\bottomrule
\end{tabular}
\end{table}

For completeness, Table~\ref{tab:results-full} reports the full
metric set including MRR and Precision@10.

\begin{table}[h]
\centering
\caption{Complete metric set, including MRR@1/5/10/20/50 and
Precision@10.}
\label{tab:results-full}
\small
\begin{tabular}{l r r r r r}
\toprule
\textbf{Dataset} & \textbf{P@10} & \textbf{MRR@10} & \textbf{MRR@50}
  & \textbf{nDCG@50} & \textbf{R@50} \\
\midrule
CLERC              & \num{0.0270} & \num{0.0444} & \num{0.0494}
  & \num{0.0870} & \num{0.1873} \\
CUAD               & \num{0.0032} & \num{0.0102} & \num{0.0148}
  & \num{0.0392} & \num{0.1458} \\
GreekBarBench      & \num{0.0261} & \num{0.0828} & \num{0.1007}
  & \num{0.2060} & \num{0.6338} \\
LaRA               & \num{0.0886} & \num{0.7697} & \num{0.7748}
  & \num{0.8197} & \num{0.9811} \\
\midrule
\textbf{Macro}     & \num{0.0362} & \num{0.2268} & \num{0.2349}
  & \num{0.2880} & \num{0.4870} \\
\bottomrule
\end{tabular}
\end{table}

\subsection{Primary metric}
The primary metric of the study,
$\mathrm{macro\_document\_Recall@10}$, is \num{0.3144}. This is the
value against which LongRAG will be compared to the other
architectures of the program under the same evaluator, query set,
and ground-truth projection.

\subsection{Qualitative shape of the results}
\begin{itemize}[leftmargin=1.5em,itemsep=0.15em]
  \item \textbf{LaRA dominates.} With $131$ source documents and
    \num{2251} long units, LaRA produces the highest scores by a
    large margin. Its corpus has a very small number of very long
    source documents, so almost every long unit is a \emph{part of}
    a distinct source document, and the retrievable granularity is
    fine at the document level.
  \item \textbf{GreekBarBench is intermediate.} With $284$ documents
    and \num{1437} units, GreekBarBench exhibits a moderate
    signal.
  \item \textbf{CLERC is weak.} With \num{2756} documents and
    \emph{only} \num{256} long units, the grouping policy packs
    roughly eleven documents into each unit. A single relevant
    passage is therefore diluted by ten unrelated passages when
    the unit is expanded into its documents. The document-level
    ranking is coarse and recall at low cutoffs is very low.
  \item \textbf{CUAD is weakest.} With $510$ contracts but
    \num{1725} long units (most of them produced by oversize
    windowing), CUAD combines two unfavourable properties: the
    relevant evidence for a query is often a \emph{single clause}
    inside a long contract, and the surrounding legal boilerplate
    dominates the dense vector. Combined with the exclusion of
    \num{68\%} of queries as zero-relevance, the residual signal
    is small.
\end{itemize}

% =====================================================================
\section{Discussion}\label{sec:discussion}

\subsection{Document-level granularity is the dominant factor}
The results are best understood in terms of a single quantity: the
average number of source documents per long unit.

\begin{center}
\small
\begin{tabular}{l r r r r}
\toprule
\textbf{Dataset} & \textbf{Docs} & \textbf{Units}
  & \textbf{Docs / unit} & \textbf{R@10} \\
\midrule
CLERC              & \num{2756} & \num{256}   & \num{10.8} & \num{0.0786} \\
CUAD               & \num{510}  & \num{1725}  & \num{0.30} & \num{0.0324} \\
GreekBarBench      & \num{284}  & \num{1437}  & \num{0.20} & \num{0.2606} \\
LaRA               & \num{131}  & \num{2251}  & \num{0.06} & \num{0.8861} \\
\bottomrule
\end{tabular}
\end{center}

The pattern is unambiguous. When the architecture's grouping policy
packs many documents per unit, the document-level expansion step
inserts irrelevant documents before relevant ones, and Recall@10
collapses. When the grouping policy produces one unit per document
fragment (as it does for LaRA), the same expansion step inserts
almost no irrelevant documents, and Recall@10 is high.

This is not a defect of the implementation. It is the
\emph{observable consequence} of the paper's design choice. LongRAG
was designed and evaluated on Wikipedia, where documents are small
and hyperlink-adjacent, and where the relevant evidence for a
multi-hop query genuinely spans several adjacent documents. Our
canonical corpora are \emph{not} Wikipedia. In corpora with
heterogeneous document sizes and without hyperlink adjacency, the
grouping policy degrades document-level retrieval quality. We
report this as a finding.

\subsection{The mid-execution fidelity correction}
We performed a mid-execution correction to the grouping policy
(\S\ref{sec:fidelity-fix}). We wish to emphasize three points about
this correction, because they matter for the interpretability of
the results.

\begin{enumerate}[leftmargin=1.5em,itemsep=0.15em]
  \item \textbf{The correction increased, not decreased, the
    fidelity of the study.} The prior policy
    (\texttt{single\_document\_if\_within\_budget}) would have
    produced $68218$ long units, none of them longer than the
    token budget of the source document. That is not LongRAG; it is
    dense passage retrieval with a different name. The corrected
    policy produces $5669$ units and is faithful to the paper's
    grouping description.
  \item \textbf{The corrected results are worse on document-level
    metrics.} This is expected: coarser retrieval granularity is
    harder to score at document level, especially when the
    relevant evidence for a query is a single sentence. We report
    the corrected numbers, not the prior ones.
  \item \textbf{The correction is fully traceable.} Both
    contract versions remain in git history; the closure registry
    records both; and the present report describes the change
    explicitly.
\end{enumerate}

\subsection{Comparison with the original paper}
We do not compare our macro numbers against the paper's headline
numbers. The comparison would be meaningless for several independent
reasons:

\begin{itemize}[leftmargin=1.5em,itemsep=0.15em]
  \item The paper evaluates long-context \emph{QA}, not retrieval.
  \item The paper uses a proprietary reader and an LLM-based
    answer scorer, both of which are disabled here.
  \item The paper's corpus is Wikipedia; ours is a heterogeneous
    canonical corpus.
  \item The paper's metric is answer accuracy, not document-level
    Recall@10.
\end{itemize}

What we do compare is \emph{relative} behaviour across the four
canonical datasets under the frozen protocol. That relative
behaviour is stable, well-behaved, and interpretable, and will be
directly comparable to the other architectures in the program.

% =====================================================================
\section{Limitations and Threats to Validity}\label{sec:limitations}

\subsection{Dataset scope}
The study executes four of the seven canonical datasets. The
excluded datasets are:

\begin{itemize}[leftmargin=1.5em,itemsep=0.15em]
  \item \textbf{RULER} (\num{4000} docs, mean $\approx \SI{8}{k}$
    tokens).
  \item \textbf{Loong} (\num{695} docs, mean $\approx \SI{300}{k}$
    tokens).
  \item \textbf{LongBench-v2} (\num{503} docs, mean $\approx
    \SI{5}{M}$ tokens).
\end{itemize}

These datasets were excluded by explicit scoping decision. The
motivation is documented in the closure manifest: their canonical
documents are an order of magnitude longer than the LongRAG unit
budget, so after oversize-windowing they would produce very large
unit corpora, and computing embeddings on them would consume a
disproportionate share of the study's GPU budget without adding
signal to the architectural comparison. The remaining four datasets
span small-document, medium-document, and large-document regimes and
preserve the architectural contrasts of interest.

This is a real limitation. The reported macro averages are
macros over four datasets, not seven, and should be interpreted
with the corresponding scope caveat.

\subsection{Controlled adaptation, not faithful reproduction}
The study is a controlled adaptation. We preserve the architecture's
defining retrieval logic; we replace the embedder, storage, and
downstream reader. Every substitution is documented. We do not claim
that our results equal the paper's results, and we do not claim
that BGE-M3 is LongRAG's native embedder. The controlled adaptation
is applied \emph{identically} across all architectures of the
program, so cross-architecture comparisons remain valid.

\subsection{Query embedding precision}
Query embeddings are FP16. This matches the document embedding
precision and is consistent with the primary execution profile in
the frozen embedding contract. A bounded reference profile (FP32 CPU
on a small stratified subset) is available in the frozen contract
for future precision-equivalence analysis but was not executed in
this run. Precision-equivalence is therefore not claimed.

\subsection{First-occurrence document expansion}
Document-level ranking is constructed by
\texttt{first\_occurrence\_by\_long\_unit\_rank}: when a document
appears in multiple retrieved long units, only its first occurrence
is scored. This is the reference policy in the adapter. It is not
the only reasonable policy: \emph{aggregate\_max} or
\emph{aggregate\_rrf} could be used to combine multiple occurrences.
We adopt the reference policy exactly, because changing it would
change the architecture.

\subsection{CPU query encoding false start}
An initial attempt to encode all \num{24520} query embeddings on CPU
was aborted after approximately \SI{45}{} minutes of wall time at an
observed throughput of $\approx \num{0.4}$ queries per second. The
computation was moved to GPU, where the same workload completes in
under five minutes. The CPU run produced no partial outputs; the
subsequent GPU run superseded it entirely. This is a processing
detail, not a scientific limitation, but we report it for full
transparency.

\subsection{Colab session lifecycle}
During the L3 embedding run, a Google Colab session reached its GPU
quota limit before the output bundles were persisted to durable
storage. The bundles were recovered by re-running the (deterministic)
embedding step with an incremental-save runner. Two embeddings per
dataset are therefore byte-identical in the final artifacts. This
detail is reported for transparency about the execution environment;
it does not affect the scientific content.

% =====================================================================
\section{Reproducibility}\label{sec:repro}

\subsection{Commit chain}
The pipeline is executed over a sequence of commits. Each commit
implements one stage and is accompanied by the tests that guard it.

\begin{center}
\small
\begin{tabular}{l l}
\toprule
\textbf{Commit} & \textbf{Stage} \\
\midrule
\texttt{5a384be} & L1 architecture and dataset-envelope freeze \\
\texttt{5a283af} & L1b fidelity correction (adjacent grouping) \\
\texttt{59b01f9} & L2 unit materialization \\
\texttt{6503732} & L3-prep Colab bundle preparation \\
\texttt{aee8800} & L4 Qdrant storage and reopen verification \\
\texttt{cd7e731} & L5b retrieval from pre-computed query embeddings \\
\texttt{0920ce5} & L6 evaluation \\
\texttt{f5ea8e6} & L7 closure \\
\bottomrule
\end{tabular}
\end{center}

\subsection{Checksum closure}
Every stage produces a manifest whose checksums cover all of its
outputs. The closure manifest
(\texttt{longrag\_full\_track\_closure\_manifest.yaml}) reports the
SHA-256 of each stage manifest together with the git revision
recorded by the pipeline.

\subsection{Resume-safe execution}
Every stage script is resume-safe. If the stage manifest is present
and marked \textsc{Succeeded}, the stage is skipped unless
\texttt{--force} is passed. This makes the entire pipeline robust
against interruption, and permits partial re-execution after
faulty stages have been corrected.

\subsection{Atomic writes}
Every output is written through a temporary file and an atomic
rename, so that no partial artifacts can be observed by downstream
stages. This applies uniformly to JSON, JSONL, and YAML outputs.

\subsection{Deterministic serialization}
JSON and JSONL outputs are serialized with \texttt{sort\_keys=True}
and \texttt{separators=(",", ":")}; JSON manifests use
\texttt{indent=2} and \texttt{sort\_keys=True}. YAML outputs use
\texttt{safe\_dump} with \texttt{sort\_keys=True} and
\texttt{allow\_unicode=True}. These settings guarantee byte-identical
output across independent runs of the same code with the same inputs.

\subsection{Available artifacts}
The following table lists the artifacts produced by the study. Items
marked \emph{in git} are committed; items marked \emph{gitignored}
are retained on disk only, per the repository's policy on dataset
files.

\begin{center}
\small
\begin{tabular}{l l l}
\toprule
\textbf{Artifact} & \textbf{Location} & \textbf{Status} \\
\midrule
Contracts & \texttt{baseline\_audit/longrag\_full\_track\_*.yaml} & in git \\
Pipeline scripts & \texttt{scripts/run\_longrag\_*.py} & in git \\
Colab bundles (input/output) & \texttt{outputs/external\_longrag/\_colab\_v1/} & gitignored \\
Qdrant instance & \texttt{outputs/external\_longrag/\_qdrant/} & gitignored \\
Query embeddings & \texttt{outputs/external\_longrag/\_query\_embeddings/} & gitignored \\
Per-stage artifacts & \texttt{outputs/external\_longrag/<\dataset>/*-v1/} & gitignored \\
Closure report (markdown) & \texttt{docs/longrag\_full\_track\_closure.md} & in git \\
Closure report (\LaTeX) & \texttt{reports/longrag\_full\_track\_closure\_report.tex} & gitignored \\
\bottomrule
\end{tabular}
\end{center}

% =====================================================================
\section{Conclusion}

We have executed, for the first time in this program, a complete
seven-stage retrieval-only evaluation of LongRAG under a shared
controlled protocol. The pipeline is fully reproducible, checksum-
closed, and resume-safe. The primary result is
$\mathrm{macro\_document\_Recall@10} = \num{0.3144}$ on four
heterogeneous canonical datasets.

The most important scientific observation from this run is that
LongRAG's document-level retrieval quality is highly sensitive to
the ratio of source documents to long units. In corpora whose
average document length is far below the architecture's unit budget,
the grouping policy packs many documents per unit, and document-level
ranking degrades. In corpora whose average document length is far
above the unit budget, the oversize-windowing path produces
one-unit-per-document-fragment and document-level ranking is
excellent. This observation is intrinsic to the architecture and
will be directly relevant to the interpretation of the remaining
architectures in the program.

The architecture is declared closed at commit
\texttt{f5ea8e6} with status
\textsc{LongRAG\_Full\_Pipeline\_Closed\_With\_Classified\_Limitations}.

% =====================================================================
\appendix

\section{Appendix A: Frozen Configuration Excerpts}

\subsection{A.1 --- Architecture freeze contract (v0.3.0)}
\begin{lstlisting}[language=yaml]
architecture:
  architecture_id: longrag
  fidelity_class: CONTROLLED_ADAPTATION
  controlled_track_eligible: true
  native_track_eligible: false
  source_architecture_commit: ba92f95

architecture_defining_components:
  - long_document_grouping
  - dense_only_retrieval

long_unit_semantics:
  construction_policy: adjacent_document_segments
  target_token_size: 4096
  token_counter: pinned_bge_m3_tokenizer
  segment_source: one_canonical_document_per_segment
  derived_parent_id: "<doc_id>:p0"
  derived_chunk_id:   "seg-<doc_id>"
  oversize_handling: split_into_target_size_windows_via_pinned_bge_m3_tokenizer
  tokenizer_id: BAAI/bge-m3
  tokenizer_revision: 5617a9f61b028005a4858fdac845db406aefb181
\end{lstlisting}

\subsection{A.2 --- Storage equivalence report (excerpt)}
\begin{lstlisting}[language=yaml]
candidate_backend: QDRANT_LOCAL_PERSISTENCE
reference_backend: EXACT_IN_MEMORY_DENSE
comparison:
  status: PASSED
conclusion:
  classification: SEMANTICALLY_NEUTRAL_FOR_CONTROLLED_TRACK
  substitution_allowed: true
results:
  top_1_overlap_mean: 1.0
  top_5_overlap_mean: 1.0
  top_10_overlap_mean: 1.0
  maximum_absolute_score_difference: 1.78e-15
  mean_absolute_score_difference: 6.29e-16
limitations:
  - Validated on the frozen synthetic acceptance fixture only.
  - Does not prove equivalence for approximate large-corpus settings.
  - Real-dataset index parameters require a separate contract.
\end{lstlisting}

\section{Appendix B: Cross-Dataset Registry Excerpts}

The full cross-dataset registry is available at
\texttt{baseline\_audit/longrag\_full\_track\_cross\_dataset\_registry.yaml}.
The essential per-dataset state is summarized here.

\begin{center}
\small
\begin{tabular}{l r r r r}
\toprule
\textbf{Dataset} & \textbf{Units} & \textbf{Emb. records}
  & \textbf{Queries} & \textbf{Evaluated} \\
\midrule
CLERC              & \num{256}   & \num{256}   & \num{1000}  & \num{1000} \\
CUAD               & \num{1725}  & \num{1725}  & \num{20910} & \num{6702} \\
GreekBarBench      & \num{1437}  & \num{1437}  & \num{284}   & \num{284} \\
LaRA               & \num{2251}  & \num{2251}  & \num{2326}  & \num{2326} \\
\bottomrule
\end{tabular}
\end{center}

\section{Appendix C: Full Metric Set}

\begin{table}[h]
\centering
\caption{Full per-dataset metric set.}
\small
\begin{tabular}{l r r r r r r r r r}
\toprule
& \multicolumn{3}{c}{\textbf{Hit}} & \multicolumn{3}{c}{\textbf{Recall}}
& \multicolumn{3}{c}{\textbf{MRR}} \\
\cmidrule(lr){2-4}\cmidrule(lr){5-7}\cmidrule(lr){8-10}
\textbf{Dataset} & @1 & @10 & @50 & @1 & @10 & @50 & @1 & @10 & @50 \\
\midrule
CLERC             & \num{0.0310} & \num{0.0940} & \num{0.2110}
                  & \num{0.0099} & \num{0.0786} & \num{0.1873}
                  & \num{0.0310} & \num{0.0444} & \num{0.0494} \\
CUAD              & \num{0.0039} & \num{0.0324} & \num{0.1458}
                  & \num{0.0039} & \num{0.0324} & \num{0.1458}
                  & \num{0.0039} & \num{0.0102} & \num{0.0148} \\
GreekBarBench     & \num{0.0352} & \num{0.2606} & \num{0.6338}
                  & \num{0.0352} & \num{0.2606} & \num{0.6338}
                  & \num{0.0352} & \num{0.0828} & \num{0.1007} \\
LaRA              & \num{0.7128} & \num{0.8861} & \num{0.9811}
                  & \num{0.7128} & \num{0.8861} & \num{0.9811}
                  & \num{0.7128} & \num{0.7697} & \num{0.7748} \\
\midrule
\textbf{Macro}    & \num{0.1957} & \num{0.3183} & \num{0.4929}
                  & \num{0.1905} & \num{0.3144} & \num{0.4870}
                  & \num{0.1957} & \num{0.2268} & \num{0.2349} \\
\bottomrule
\end{tabular}
\end{table}

\begin{table}[h]
\centering
\caption{Full per-dataset metric set (continued): nDCG and
Precision@10.}
\small
\begin{tabular}{l r r r}
\toprule
\textbf{Dataset} & \textbf{nDCG@10} & \textbf{nDCG@50}
  & \textbf{Precision@10} \\
\midrule
CLERC             & \num{0.0543} & \num{0.0870} & \num{0.0270} \\
CUAD              & \num{0.0153} & \num{0.0392} & \num{0.0032} \\
GreekBarBench     & \num{0.1233} & \num{0.2060} & \num{0.0261} \\
LaRA              & \num{0.7977} & \num{0.8197} & \num{0.0886} \\
\midrule
\textbf{Macro}    & \num{0.2476} & \num{0.2880} & \num{0.0362} \\
\bottomrule
\end{tabular}
\end{table}

\section{Appendix D: Timing Summary}

\begin{center}
\small
\begin{tabular}{l r}
\toprule
\textbf{Stage} & \textbf{Wall time} \\
\midrule
L1 freeze (4 datasets)                & \num{19.8}\,\si{s} \\
L2 unit materialization (4 datasets)  & \num{180.1}\,\si{s} \\
L3 embedding (Colab T4, 4 datasets)   & $\approx$ \num{1600}\,\si{s} \\
L4 Qdrant storage + reopen            & \num{16.0}\,\si{s} \\
L5 query encoding (Colab T4)          & $\approx$ \num{240}\,\si{s} \\
L5 retrieval (Qdrant + expansion)     & \num{342.0}\,\si{s} \\
L6 evaluation                         & \num{5.0}\,\si{s} \\
L7 closure                            & \num{0.1}\,\si{s} \\
\bottomrule
\end{tabular}
\end{center}

\section{Appendix E: Zero-Relevance Query Accounting}

For CUAD, \num{14208} of \num{20910} queries have no relevant
documents in the canonical ground truth. Following the frozen
evaluator contract, these queries are excluded from all metric
aggregations. For every other executed dataset, all queries have at
least one relevant document and none are excluded.

The exclusion is total: it affects Hit, Recall, Precision, MRR, and
nDCG equally. The count is reported in every per-dataset result so
that downstream analyses can reweight if needed.

\section{Appendix F: Files Produced}

For completeness, the following table lists every artifact
produced by the pipeline. Sizes are reported in bytes as recorded in
the L7 closure manifest.

\begin{center}
\small
\begin{tabular}{l r}
\toprule
\textbf{Artifact} & \textbf{Bytes} \\
\midrule
\texttt{outputs/external\_longrag/clerk/units-v1/units.jsonl}             & \num{4280790} \\
\texttt{outputs/external\_longrag/cuad/units-v1/units.jsonl}              & \num{26656099} \\
\texttt{outputs/external\_longrag/greek\_bar\_bench/units-v1/units.jsonl} & \num{35116432} \\
\texttt{outputs/external\_longrag/lara/units-v1/units.jsonl}              & \num{34637258} \\
\texttt{outputs/external\_longrag/clerk/embeddings-v1/embedding\_records.jsonl}             & \num{5527272} \\
\texttt{outputs/external\_longrag/cuad/embeddings-v1/embedding\_records.jsonl}              & \num{35931570} \\
\texttt{outputs/external\_longrag/greek\_bar\_bench/embeddings-v1/embedding\_records.jsonl} & \num{30024682} \\
\texttt{outputs/external\_longrag/lara/embeddings-v1/embedding\_records.jsonl}              & \num{46875519} \\
\texttt{outputs/external\_longrag/\_qdrant/}                                & $\approx$ \num{70238123} \\
\bottomrule
\end{tabular}
\end{center}

\end{document}